% Generated by tables/make_methods-compared.py from data/derived/benchmark/methods-compared.json; do not edit.
\subsubsection{Default prediction on a retail book}
\label{sec:results-methods-compared-q5}

This question is answered on case C01. A locked stratified holdout of 30\%, with repeated stratified five-fold cross-validation inside the training part; every rung of one document is scored on the same held-out rows. The two documents hold different data and are never ranked together. The calibration slope is not in any committed document; the expected calibration error and the Spiegelhalter p-value are. The run ``Taiwan credit-card holdout'' reads as follows: 9,000 held-out rows; its random seed is 42. The run ``German credit twin'' reads as follows: 300 held-out rows; the only document that runs TabPFN; its random seed is 42. Table~\ref{tab:bm-q5-methods} lists the nine methods the document compares, two long-established, five recent published and two reference, with the section of this report that describes each. Tables~\ref{tab:bm-q5-1}, \ref{tab:bm-q5-2} give the values, one table for each group of metrics that share a run and a layout.

{\footnotesize\setlength{\tabcolsep}{3pt}
% bm-table-methods Q5
\begin{longtable}{>{\raggedright\arraybackslash}p{0.38\textwidth}>{\raggedright\arraybackslash}p{0.17\textwidth}>{\raggedright\arraybackslash}p{0.27\textwidth}>{\raggedright\arraybackslash}p{0.10\textwidth}}
\caption{The methods compared on the question Default prediction on a retail book: the standing of each method, the source and the year the document records for it, and the section of this report that describes it; a dash marks a method with no source recorded or with no section of its own.}\label{tab:bm-q5-methods}\\
\toprule Method & Standing & Source & Section \\ \midrule \endfirsthead
\toprule Method & Standing & Source & Section \\ \midrule \endhead
Weight-of-evidence logistic scorecard & Long-established & -- & \ref{sec:m1} \\
L1-penalised logistic regression on weight-of-evidence & Long-established & -- & \ref{sec:m2} \\
Penalised logistic tree regression & Recent published & -- & \ref{sec:m3} \\
Explainable boosting machine & Recent published & -- & \ref{sec:m4} \\
Monotone gradient boosting (LightGBM) with isotonic calibration & Recent published & -- & \ref{sec:m5} \\
Monotone gradient boosting (XGBoost), second engine & Recent published & -- & \ref{sec:m5} \\
TabPFN & Recent published & -- & \ref{sec:m6} \\
Constant PD & Reference & -- & -- \\
One-variable PD & Reference & -- & -- \\
\bottomrule
\end{longtable}
}

{\footnotesize\setlength{\tabcolsep}{3pt}
% bm-table Q5 normal cols=tw_auc/holdout;tw_brier/holdout;tw_ece/holdout;tw_sph/holdout
\begin{longtable}{>{\raggedright\arraybackslash}p{0.25\textwidth}>{\raggedleft\arraybackslash}p{0.14\textwidth}>{\raggedleft\arraybackslash}p{0.14\textwidth}>{\raggedleft\arraybackslash}p{0.14\textwidth}>{\raggedleft\arraybackslash}p{0.14\textwidth}}
\caption{Default prediction on a retail book: Area under the ROC curve (AUC; higher is better); Brier score (Brier; lower is better); Expected calibration error (error; lower is better); Spiegelhalter calibration p-value (evidence against calibration, not a score) (p-value; context, not ranked), for each method, on the run Taiwan credit-card holdout. The best value of a column by the declared rule is in bold, ties all in bold; a dash marks a cell without a value, and the notes below the table state why.}\label{tab:bm-q5-1}\\
\toprule Method & Area under the ROC curve & Brier score & Expected calibration error & Spiegelhalter calibration p-value (evidence against calibration, not a score) \\ \midrule \endfirsthead
\toprule Method & Area under the ROC curve & Brier score & Expected calibration error & Spiegelhalter calibration p-value (evidence against calibration, not a score) \\ \midrule \endhead
\multicolumn{5}{l}{\emph{Long-established}} \\
Weight-of-evidence logistic scorecard & 0.759 & 0.1393 & 0.016 & 0.022 \\ % bm-row P1-scorecard
L1-penalised logistic regression on weight-of-evidence & 0.759 & 0.1393 & 0.015 & 0.024 \\ % bm-row P2a-l1
\multicolumn{5}{l}{\emph{Recent published}} \\
Penalised logistic tree regression & 0.775 & 0.1367 & \textbf{0.011} & 0.084 \\ % bm-row P2b-pltr
Explainable boosting machine & 0.775 & 0.1367 & 0.012 & 0.231 \\ % bm-row P3-ebm
Monotone gradient boosting (LightGBM) with isotonic calibration & \textbf{0.777} & 0.1361 & 0.013 & 0.115 \\ % bm-row P4-lightgbm
Monotone gradient boosting (XGBoost), second engine & 0.775 & \textbf{0.1359} & 0.012 & 0.240 \\ % bm-row P4-xgboost
TabPFN & -- & -- & -- & -- \\ % bm-row P5-tabpfn
\multicolumn{5}{l}{\emph{Reference}} \\
Constant PD & 0.500 & 0.1723 & 0.017 & 0.497 \\ % bm-row P0-constant
One-variable PD & 0.692 & 0.1431 & 0.014 & 0.110 \\ % bm-row P0-single
\bottomrule
\end{longtable}
}

A dash in Table~\ref{tab:bm-q5-1} is a cell the product does not compute, for the stated reason: TabPFN runs only on the small cases (the German document), never on the 9,000-row holdout (TabPFN).

{\footnotesize\setlength{\tabcolsep}{3pt}
% bm-table Q5 normal cols=de_auc/german-twin;de_brier/german-twin;de_ece/german-twin;de_sph/german-twin
\begin{longtable}{>{\raggedright\arraybackslash}p{0.25\textwidth}>{\raggedleft\arraybackslash}p{0.14\textwidth}>{\raggedleft\arraybackslash}p{0.14\textwidth}>{\raggedleft\arraybackslash}p{0.14\textwidth}>{\raggedleft\arraybackslash}p{0.14\textwidth}}
\caption{Default prediction on a retail book: Area under the ROC curve (AUC; higher is better); Brier score (Brier; lower is better); Expected calibration error (error; lower is better); Spiegelhalter calibration p-value (evidence against calibration, not a score) (p-value; context, not ranked), for each method, on the run German credit twin. The best value of a column by the declared rule is in bold, ties all in bold; a dash marks a cell without a value, and the notes below the table state why.}\label{tab:bm-q5-2}\\
\toprule Method & Area under the ROC curve & Brier score & Expected calibration error & Spiegelhalter calibration p-value (evidence against calibration, not a score) \\ \midrule \endfirsthead
\toprule Method & Area under the ROC curve & Brier score & Expected calibration error & Spiegelhalter calibration p-value (evidence against calibration, not a score) \\ \midrule \endhead
\multicolumn{5}{l}{\emph{Long-established}} \\
Weight-of-evidence logistic scorecard & 0.788 & 0.162 & 0.045 & 0.407 \\ % bm-row P1-scorecard
L1-penalised logistic regression on weight-of-evidence & \textbf{0.791} & \textbf{0.161} & 0.052 & 0.617 \\ % bm-row P2a-l1
\multicolumn{5}{l}{\emph{Recent published}} \\
Penalised logistic tree regression & 0.786 & 0.166 & 0.048 & 0.371 \\ % bm-row P2b-pltr
Explainable boosting machine & 0.787 & 0.167 & 0.058 & 0.0074 \\ % bm-row P3-ebm
Monotone gradient boosting (LightGBM) with isotonic calibration & 0.778 & 0.173 & 0.076 & \ensuremath{2.2\times10^{-7}} \\ % bm-row P4-lightgbm
Monotone gradient boosting (XGBoost), second engine & -- & -- & -- & -- \\ % bm-row P4-xgboost
TabPFN & 0.777 & 0.167 & \textbf{0.039} & 0.0079 \\ % bm-row P5-tabpfn
\multicolumn{5}{l}{\emph{Reference}} \\
Constant PD & 0.500 & 0.210 & 0.047 & 0.500 \\ % bm-row P0-constant
One-variable PD & 0.723 & 0.182 & 0.074 & 0.713 \\ % bm-row P0-single
\bottomrule
\end{longtable}
}

A dash in Table~\ref{tab:bm-q5-2} is a cell the product does not compute, for the stated reason: the German document runs one gradient-boosting engine only (Monotone gradient boosting (XGBoost), second engine).

The document reads the tables of this question by the declared rule of each metric. Area under the ROC curve (higher is better): best in Monotone gradient boosting (LightGBM) with isotonic calibration 1 of 1. Brier score (lower is better): best in Monotone gradient boosting (XGBoost), second engine 1 of 1. Expected calibration error (lower is better): best in Penalised logistic tree regression 1 of 1. Area under the ROC curve (higher is better): best in L1-penalised logistic regression on weight-of-evidence 1 of 1. Brier score (lower is better): best in L1-penalised logistic regression on weight-of-evidence 1 of 1. Expected calibration error (lower is better): best in TabPFN 1 of 1. A null cell is never ranked; each states why it is null.
